\section{Two-prime support and rational-return rigidity}
\label{sec:uniform-analytic-continuation}

We give a complete elementary bound on one moving-support subclass of
abc triples, and then examine a proposed passage from local logarithmic
arrows to rational $S$-unit solutions. These are different conclusions:
the first is a special-case height estimate, while the second restricts
the number of outputs of a specified construction. Throughout this
section, $r=\rad(abc)$ and $R=\log r$, as in \eqref{eq:abc}.

\subsection{The complete two-prime subclass}

\begin{theorem}[A sharp bound on two-prime support]
\label{thm:uac-two-prime}
Let $a,b,c$ be positive integers with $a+b=c$ and $\gcd(a,b)=1$.
If $\omega(abc)\le2$, then
\begin{equation}\label{eq:uac-two-prime-bound}
                         c\le\frac32\rad(abc).
\end{equation}
Equality holds precisely for $(a,b,c)=(1,8,9)$ or $(8,1,9)$.
Every other triple in this class satisfies $c\le\rad(abc)$.
\end{theorem}

\begin{proof}
The entries are pairwise coprime. Exactly one is even: the addends
cannot both be even, and their parities determine that of their sum.
If all three entries exceeded $1$, each would contribute a different
prime divisor, contradicting the support bound. An addend is therefore
$1$. The case $(1,1,2)$ has $c=r=2$. Otherwise, after exchanging the
addends if necessary, write the triple as $(1,N,N+1)$, with $N\ge2$.
Both consecutive integers have a prime factor, and their supports are
disjoint. Thus, for an odd prime $q$ and positive integers $u,v$,
\begin{equation}\label{eq:uac-consecutive-powers}
                         \{N,N+1\}=\{2^u,q^v\}.
\end{equation}

Suppose first that $2^u=q^v+1$. If $v$ is even, then
$q^v\equiv1\pmod4$, so $2^u\equiv2\pmod4$ and $u=1$,
contrary to $q^v+1>2$. If $v\ge3$ is odd, the factorization
\[
 q^v+1=(q+1)(q^{v-1}-q^{v-2}+\cdots-q+1)
\]
has an odd second factor greater than $1$. Indeed, its $v$ terms
are odd modulo $2$, and its value is at least $q^2-q+1>1$.
Such a factor cannot divide a power of $2$. This ordering therefore
requires $v=1$.

Now suppose that $q^v=2^u+1$. An odd $v\ge3$ is excluded by
\[
 q^v-1=(q-1)(1+q+\cdots+q^{v-1}),
\]
whose second factor is again odd and greater than $1$.
If $v$ is even, set $Q=q^{v/2}\ge3$. Both factors in
\[
                         (Q-1)(Q+1)=2^u
\]
are positive powers of $2$. Write $Q-1=2^s$ and $Q+1=2^t$,
with $1\le s<t$. Subtraction gives
$2^s(2^{t-s}-1)=2$, hence $s=1$, $t=2$ and $Q=3$.
Since $q$ is prime, this forces $q=3$, $v=2$ and $u=3$.
The resulting pair is $\{8,9\}$.

In every remaining case $v=1$, the radical is $r=2q$, and the
larger endpoint is $q$ or $q+1$, both at most $2q$.
For $(1,8,9)$ and its exchanged triple, $c=9$ and $r=6$.
These observations prove both the bound and its equality statement.
\end{proof}

The estimate is uniform as the odd prime varies. For every $\eps>0$,
the entire subclass in Theorem~\ref{thm:uac-two-prime} satisfies
\[
 \log c\le R+\log(3/2)
        \le(1+\eps)R+\log(3/2),
\]
with no constant depending on the prime support. The proof is elementary
and uses no classification theorem for general consecutive perfect
powers. No priority claim is made for this classical special case, and
no estimate for $\omega(abc)\ge3$ follows from it.

\subsection{Nonzero rational return under the actual trace constraint}

Let $E,E'$ be finite extensions of $\Q_p$. We retain the
coefficient-compatible arrow of Proposition~\ref{prop:trace-transport}:
a continuous isomorphism $\alpha:G_{E'}\simeq G_E$ together with
an integral equivariant coefficient isomorphism
\[
 \beta:\alpha^*\Z_p(1)_E\simeq\Z_p(1)_{E'}.
\]
Its induced logarithmic map $F:E\longrightarrow E'$ is $\Q_p$-linear
and satisfies, for some $\kappa\in\Z_p^\times$,
\begin{equation}\label{eq:uac-trace-covariance}
 \operatorname{Tr}_{E'/\Q_p}(Fz)
                 =\kappa\,\operatorname{Tr}_{E/\Q_p}(z)
                 \qquad(z\in E).
\end{equation}
This is the previously proved identity \eqref{eq:trace-transport}.
Its unit multiplier comes from the integral transport of
$H^2(G_E,\Z_p(1))\cong\Z_p$, together with the reciprocity
cup-product formula
\cite[Chapter~III, Proposition~3.6 and \S4]{MilneCFT}.
Cyclotomic character, inertia and absolute degree are recovered from
the local Galois group
\cite[Propositions~1.1--1.2 and Corollary~1.3]{MochLocal}.
In particular the degrees of $E$ and $E'$ agree for these arrows.
None of these assertions requires the arrow to be induced by a field
isomorphism.

\begin{lemma}[Rational return and the scalar line]
\label{lem:uac-rational-return}
Suppose $[E:\Q_p]=[E':\Q_p]=d$ and a $\Q_p$-linear map $F$
satisfies \eqref{eq:uac-trace-covariance}.
If $z\in\Q_p\subset E$ and $F(z)=z'\in\Q_p\subset E'$, then
\begin{equation}\label{eq:uac-rational-scalar}
                              z'=\kappa z.
\end{equation}
For nonzero $z$, the output is nonzero and $v_p(z')=v_p(z)$.
Moreover, one such nonzero return implies
\[
                         F(t)=\kappa t\qquad(t\in\Q_p).
\]
\end{lemma}

\begin{proof}
The traces of the two rational elements are $dz$ and $dz'$.
Equation~\eqref{eq:uac-trace-covariance} gives $dz'=\kappa dz$.
The positive integer $d$ is nonzero in $\Q_p$, so cancellation
gives \eqref{eq:uac-rational-scalar}, even when $p\mid d$.
Its valuation assertion follows from $v_p(\kappa)=0$.
If $z\ne0$, linearity gives
$F(1)=z^{-1}F(z)=\kappa$, and then $F(t)=tF(1)=\kappa t$.
\end{proof}

Equal degree is an explicit hypothesis in this linear-algebra statement.
With unequal degrees $d,d'$, the same calculation would instead yield
$z'=\kappa(d/d')z$. Uniform division of both logarithmic coordinates
by $p$ leaves the displayed trace identity unchanged; independent
rescalings by different powers of $p$ do not. The lemma does not assert
that every local Galois arrow preserves the rational line. It describes
what follows if a specified point returns to that line.

\begin{lemma}[Depth of an odd-prime principal-unit logarithm]
\label{lem:uac-log-depth}
For an odd prime $p$ and $0\ne\xi\in p\Z_p$,
\begin{equation}\label{eq:uac-log-depth}
                  v_p\bigl(\log_p(1-\xi)\bigr)=v_p(\xi).
\end{equation}
\end{lemma}

\begin{proof}
Put $e=v_p(\xi)\ge1$. In the convergent series
\[
                    \log_p(1-\xi)=-\sum_{n\ge1}\frac{\xi^n}{n},
\]
the first term has valuation $e$. For $n\ge2$ one has
$v_p(n)\le n-2$: this is immediate for $n=2$, and for $n\ge3$
it follows from $v_p(n)\le\log_3 n\le n-2$.
Every subsequent term therefore has valuation
$ne-v_p(n)\ge e+1$, and these valuations tend to infinity.
The first term is the unique term of least valuation. The logarithm
is nonzero and has valuation $e$.
\end{proof}

For an odd prime $p\mid abc$, choose the following vanishing coordinate
and its companion rational $S$-unit, where $S$ is the full prime support
of $abc$:
\begin{equation}\label{eq:uac-companions}
\begin{array}{c|cc}
 &\xi_p&u_p=1-\xi_p\\ \hline
 p\mid a&a/c&b/c\\
 p\mid b&b/c&a/c\\
 p\mid c&c/a&-b/a .
\end{array}
\end{equation}
Pairwise coprimality makes every denominator a $p$-adic unit. Thus
$\xi_p\in p\Z_p$, $u_p\in1+p\Z_p$, and $v_p(\xi_p)$ is exactly
the exponent of $p$ in the indicated endpoint. In the last row
$1-c/a=-b/a$; the negative real sign does not prevent
$p$-adic congruence to $1$.

\begin{proposition}[An endpoint exponent cannot change under rational return]
\label{prop:uac-endpoint-depth}
Let $(a,b,c)$ and $(A,B,C)$ be positive primitive triples, and let
an odd prime $p$ divide the same labelled endpoint in both triples.
Form $u_p,u'_p$ using that row of \eqref{eq:uac-companions}.
Suppose an equal-degree coefficient-compatible logarithmic arrow
satisfies the pointwise equality
\begin{equation}\label{eq:uac-point-match}
                         F_p(\log_p u_p)=\log_p u'_p .
\end{equation}
Then the exponent of $p$ in the indicated endpoint is the same in
the two triples.
\end{proposition}

\begin{proof}
Both logarithms are nonzero rational $p$-adic elements by
Lemma~\ref{lem:uac-log-depth}. Lemma~\ref{lem:uac-rational-return}
equates their valuations, and \eqref{eq:uac-log-depth} identifies
those valuations with the corresponding endpoint exponents.
\end{proof}

The proposition imposes no condition at $p=2$.
The depth identity used here is not valid for every
$\xi\in2\Z_2$. Nor can \eqref{eq:uac-point-match} be replaced by
membership of the two logarithms in the same ideal or hull.

\subsection{The exact arithmetic fibre bound}

For $n>0$, write $n_{\rm odd}=n/2^{v_2(n)}$. Fix a positive primitive
triple $(a,b,c)$ and put
$(a_0,b_0,c_0)=(a_{\rm odd},b_{\rm odd},c_{\rm odd})$.
Let $\mathcal F(a,b,c)$ be the set of all positive primitive triples
$(A,B,C)$ with
\[
                    (A_{\rm odd},B_{\rm odd},C_{\rm odd})
                                      =(a_0,b_0,c_0).
\]
These conditions fix all odd-prime exponents in their labelled
endpoints, while allowing the even endpoint to move.

\begin{theorem}[Two possible outputs]
\label{thm:uac-odd-fibre}
For every positive primitive triple,
\begin{equation}\label{eq:uac-fibre-bound}
                              |\mathcal F(a,b,c)|\le2.
\end{equation}
If the position of the even endpoint is fixed to that of the seed,
the fibre contains only the seed. The bound two is attained.
\end{theorem}

\begin{proof}
Every target has exactly one even endpoint, with exponent $k\ge1$.
The three possible positions give
\begin{equation}\label{eq:uac-fibre-cases}
\begin{array}{c|c|c}
 \text{even endpoint}&(A,B,C)&\text{required equation}\\ \hline
 A&(2^ka_0,b_0,c_0)&2^ka_0+b_0=c_0\\
 B&(a_0,2^kb_0,c_0)&a_0+2^kb_0=c_0\\
 C&(a_0,b_0,2^kc_0)&a_0+b_0=2^kc_0 .
\end{array}
\end{equation}
Each row admits at most one $k$, since $2^k$ multiplies a fixed
positive integer and is strictly increasing. Either of the first
two rows implies $c_0>a_0+b_0$, whereas the third implies
$c_0<a_0+b_0$. The third row cannot coexist with either of the
first two, so at most two triples occur.

When the even position is fixed, only one row remains, and it
already contains the seed. No division by a difference of odd
parts was used. In particular, if $a_0=b_0$, coprimality forces
$a_0=b_0=1$ and the same argument applies. When all three odd
parts are $1$, the only possibility is $(1,1,2)$.

Finally, the two distinct triples
\begin{equation}\label{eq:uac-sharp-triples}
                          (4,3,7),\qquad(1,6,7)
\end{equation}
are primitive and have the same labelled odd parts $(1,3,7)$.
Thus the bound is sharp.
\end{proof}

\begin{corollary}[The specified rational-return construction]
\label{cor:uac-rational-return-fibre}
Fix a seed $(a,b,c)$ with full prime support $S$. Consider positive
primitive targets $(A,B,C)$ satisfying all of the following:
\begin{enumerate}
\item Their exact prime support is $S$.
\item Every odd $p\in S$ belongs to the same labelled endpoint as
in the seed.
\item For each odd $p\in S$, an equal-degree coefficient-compatible
local logarithmic arrow gives \eqref{eq:uac-point-match}, with the
companion units in \eqref{eq:uac-companions}.
\end{enumerate}
There are at most two targets. If the even endpoint is prescribed
as well, the only target is the seed.
\end{corollary}

\begin{proof}
Proposition~\ref{prop:uac-endpoint-depth} preserves every odd exponent
on the specified support. Primitivity makes its exponents on the other
two endpoints zero. Exact support equality excludes any new odd primes,
so all three odd parts agree. Apply Theorem~\ref{thm:uac-odd-fibre}.
Identity arrows admit the seed itself.
\end{proof}

The local arrows may depend on the prime and the target; the upper
bound requires no common global arrow. Exact support and endpoint-label
preservation are hypotheses, not consequences of trace transport.
Checking the old primes alone would not exclude new prime factors.

\begin{proposition}[Sharpness with freely chosen coefficient units]
\label{prop:uac-coefficient-sharpness}
In the class of pairs $(\alpha,\beta)$ in which an identity-Galois
arrow may be accompanied by any integral Tate-module coefficient unit,
the upper bound in Corollary~\ref{cor:uac-rational-return-fibre} is attained
by the seed and target in \eqref{eq:uac-sharp-triples}.
\end{proposition}

\begin{proof}
Both triples have support $\{2,3,7\}$; the odd prime $3$ stays on
the second endpoint and $7$ on the third. Their companion units are
\[
 (u_3,u'_3)=(4/7,1/7),\qquad
 (u_7,u'_7)=(-3/4,-6).
\]
The vanishing coordinates at $3$ are $3/7,6/7$, and those at $7$
are $7/4,7$. All have valuation one at their indicated prime.
Lemma~\ref{lem:uac-log-depth} therefore shows that
\[
                       \kappa_p=
                     \frac{\log_p u'_p}{\log_p u_p}
                        \in\Z_p^\times\qquad(p=3,7).
\]
Take $E=E'=\Q_p$ and $\alpha=\operatorname{id}_{G_{\Q_p}}$.
Multiplication by $\kappa_p$ on $\Z_p(1)$ is an integral
equivariant coefficient isomorphism. It acts as that same scalar
on unit Kummer cohomology and its logarithmic module, so it gives
exactly \eqref{eq:uac-point-match}. These are principal units at
odd primes, where logarithm is injective; no torsion ambiguity
affects the stated points. Together with the identity arrows for
the seed, this supplies two targets in the specified coefficient-pair
class.
\end{proof}

The coefficient convention in this proposition is essential.
The purely arithmetic sharpness of Theorem~\ref{thm:uac-odd-fibre}
has no Galois hypothesis. Its additional realization just proved
allows $\beta$ to be multiplied by a freely chosen element of
$\Z_p^\times$. It does not show that every resulting $\kappa_p$
is induced by a strict automorphism of $G_{\Q_p}$ when a Tate
coefficient framing has been fixed. For that possibly smaller
class, the upper bound remains valid, but this example alone
does not establish sharpness. No arrow condition at $p=2$ was imposed.

\subsection{Formal scope and the unproved uniform step}

Theorem~\ref{thm:uac-two-prime}, including its equality classification,
is formalized in the separate module
\texttt{ABCTwoPrimeSupport20260831}. Its proof starts with the actual
prime-factor support and the repository radical; the prime-power
classification is proved rather than assumed. All 22 public theorems
in that module were compiled and checked for axiom dependencies.
The module \texttt{TraceCovariantRationalReturn20260831} proves the
algebraic cancellation and scalar-line conclusions using the actual
algebra trace and actual dimensions. Trace covariance remains an
explicit hypothesis there: this does not formalize its local Galois
or Kummer derivation, or the principal-unit logarithm argument.
The separate module \texttt{ABCOddPartFibre20260831} now proves
Theorem~\ref{thm:uac-odd-fibre} for the actual natural-number odd part
and the unchanged primitive-triple type. It constructs an injection
of each entire fibre into a two-element type, proves that the fibre
is finite, and establishes cardinality exactly two for the displayed
example. This is an arithmetic fibre theorem; the logarithmic
endpoint-transfer proposition and the coefficient-pair realization
remain the mathematical proofs above, rather than formalized
Galois-orbit assertions.

The rational-return obstruction is compatible with the minimum-layer
constructions of the preceding sections. Those constructions can move
a rational input outside the rational line and then form a hull.
Returning to an actual rational $S$-unit solution is an additional
condition. Multiplication by tensor-order coefficients, addition and
ideal generation do not automatically preserve a point's trace, so
the fibre bound cannot be applied after replacing
\eqref{eq:uac-point-match} by hull membership.

No general $S$-unit method, IUT theorem, complete Ind3 operation or abc
conjecture is disproved by this obstruction. Mechanisms introducing new
prime support, changing prime labels, or using nonrational intermediate
points followed by a proved arithmetic descent remain outside its
scope. The uniform estimate \eqref{eq:abc} for $\omega(abc)\ge3$,
even for a fixed number of primes which themselves vary, is not proved
here. Neither finite exceptional sets depending on the support nor
density-one bounds supply that missing quantifier.
